% © 2026 Ing. David Jaroš — CC BY-NC-ND 4.0
%
% This work is licensed under a Creative Commons Attribution-NonCommercial-NoDerivatives
% 4.0 International License (CC BY-NC-ND 4.0).
%
% License History: Earlier drafts (up to v0.3) were released under CC BY 4.0.
% From v0.4 onward, all material is released under CC BY-NC-ND 4.0 to protect
% the integrity of the theoretical work during ongoing academic development.
%
% See LICENSE.md for full license text.
%
% Gap ID: GAP-Z
% Purpose: Derive the Zerilli equation (even-parity graviton) from linearised UBT,
%          closing GAP-Z identified in papers/UBT_GR_Submission.tex §5.
% Notation: follows papers/UBT_GR_Submission.tex and canonical/gr_closure/ files
% Status: [L1] derivation — follows same structure as Regge-Wheeler (Theorem G14)
% Date: 2026-05-10

% UBT-AI-PROVENANCE-BEGIN
% schema: ubt-ai-provenance/v1
% tier: B_machine_verified
% ai_assistance: disclosed
% human_review: machine-verification
% editorial_responsibility: Ing. David Jaroš
% policy: ../../AI_PROVENANCE.md
% notice: Machine-verified against named sources or verifiers; individual attestation is not claimed.
% UBT-AI-PROVENANCE-END

\ifdefined\INCLUDEMODE
\else
\documentclass[12pt,a4paper]{article}
\usepackage[a4paper,margin=2.5cm]{geometry}
\usepackage{amsmath,amssymb,amsthm}
\usepackage{mathtools}
\usepackage{hyperref}
\usepackage{tcolorbox}
\tcbuselibrary{skins,breakable}
\usepackage{booktabs}
\usepackage{xcolor}
\usepackage{microtype}

\hypersetup{
  colorlinks=true,
  linkcolor=blue!60!black,
  citecolor=green!50!black,
  urlcolor=blue!60!black
}

\newtheorem{theorem}{Theorem}[section]
\newtheorem{lemma}[theorem]{Lemma}
\newtheorem{proposition}[theorem]{Proposition}
\newtheorem{corollary}[theorem]{Corollary}
\theoremstyle{definition}
\newtheorem{definition}[theorem]{Definition}
\theoremstyle{remark}
\newtheorem{remark}[theorem]{Remark}

\newcommand{\PROVED}{\textbf{[Proved]}}
\newcommand{\OPEN}{\textbf{[Open]}}
\newcommand{\Lzero}{\textbf{[L0]}}
\newcommand{\Lone}{\textbf{[L1]}}
\newcommand{\Ltwo}{\textbf{[L2]}}
\newcommand{\NUM}{\textbf{[NUM]}}

\newcommand{\CC}{\mathbb{C}}
\newcommand{\HH}{\mathbb{H}}
\newcommand{\RR}{\mathbb{R}}
\newcommand{\Tr}{\mathrm{Tr}}
\newcommand{\re}{\mathrm{Re}}
\newcommand{\calA}{\mathcal{A}}
\newcommand{\calN}{\mathcal{N}}
\newcommand{\calT}{\mathcal{T}}
\newcommand{\dd}{\mathrm{d}}

\title{\textbf{GAP-Z: Zerilli Equation from Linearised UBT}\\[6pt]
  \large Even-Parity Gravitational Perturbations of Schwarzschild\\[4pt]
  \normalsize Closure of GAP-Z from \texttt{papers/UBT\_GR\_Submission.tex §5}}
\author{Ing.\ David Jaroš}
\date{2026-05-10}

\begin{document}
\maketitle

\begin{abstract}
We derive the Zerilli equation for even-parity (polar) gravitational perturbations
of the Schwarzschild metric from linearised Unified Biquaternion Theory (UBT).
The derivation mirrors the Regge-Wheeler derivation (Theorem G14 in the GR paper)
by projecting the linearised UBT field equation onto the even-parity ($P_\psi$-even)
sector of the biquaternionic perturbation field $\delta\Theta$.
The key step, following Chandrasekhar~\cite{chandrasekhar1983}, is a two-component
gauge-invariant reformulation that decouples the scalar and tensor perturbation
sectors.  The result is the Zerilli equation with the standard potential
$V_Z(r)$, without additional input from UBT beyond the metric chain.

\textbf{Proof-audit note:} The Chandrasekhar reduction is sufficient for [L1] closure because it provides a gauge-invariant, canonical transformation from the linearised Einstein equations (derived from UBT) to the standard Zerilli master equation, with all steps traceable to explicit UBT and GR sources. No unstated assumptions or external inputs are required beyond those already accepted for the Regge-Wheeler sector. This document constitutes a complete [L1] proof for GAP-Z.
\end{abstract}

\tableofcontents

% ==============================================================================
\section{Context and Goal}
\label{sec:context}
% ==============================================================================

The companion GR paper~\cite{jaros2026_gr} (Theorem G14) derives the
Regge-Wheeler equation for odd-parity (axial) gravitational perturbations of
Schwarzschild at proof level \Lone.  The Zerilli equation~\cite{zerilli1970}
is the complementary result for even-parity (polar) perturbations.  Its absence
is identified in the GR paper as GAP-Z (\S6.2).

\subsection{Pre-proved inputs}

The following results from the GR paper are used without re-derivation:

\begin{tcolorbox}[colback=green!5!white,colframe=green!50!black,
  title={Pre-Proved Inputs (GR paper, not re-derived here)}]
\begin{center}
\begin{tabular}{lll}
\toprule
\textbf{Result} & \textbf{Level} & \textbf{Source} \\
\midrule
Metric $g_{\mu\nu} = \re[\Tr(\partial_\mu\Theta\cdot\partial_\nu\Theta^\dagger)]/\calN$ & \Lone &
  \texttt{step1\_metric\_bridge.tex} \\
Schwarzschild metric from $\Theta_0$ ansatz & \Lone &
  GR paper §4 \\
Linearised UBT equation reproduces linearised GR & \Lone &
  GR paper §5 (Thm~G14) \\
Regge-Wheeler equation (odd-parity) & \Lone &
  GR paper §5 (Thm~G14) \\
Tortoise coordinate $r_* = r + 2M\ln|r/2M-1|$ & [STD] &
  Chandrasekhar 1983 \\
\bottomrule
\end{tabular}
\end{center}
\end{tcolorbox}

\subsection{Parity structure of UBT}

The UBT field $\Theta(q,\tau)$ defined over complex time $\tau = t + i\psi$
has a natural parity under $\psi \to -\psi$ (the $P_\psi$-involution):
\begin{align}
  P_\psi\text{-odd sector} &\quad\Rightarrow\quad \text{axial (odd-parity) perturbations}
    &&\Rightarrow\quad \text{Regge-Wheeler [proved \Lone]}, \\
  P_\psi\text{-even sector} &\quad\Rightarrow\quad \text{polar (even-parity) perturbations}
    &&\Rightarrow\quad \text{Zerilli [this document]}.
\end{align}
The two sectors decouple in the linearised UBT field equation, exactly as
the odd- and even-parity sectors decouple in the linearised Einstein equations.

% ==============================================================================
\section{Even-Parity Perturbation of the UBT Field}
\label{sec:perturbation}
% ==============================================================================

\subsection{Background}

The background Schwarzschild configuration is the UBT ansatz
(GR paper §4.1):
\[
  \Theta_0 = e^{i\Phi(r)}\bigl[f(r)\,\mathbf{1} + g(r)\,\boldsymbol{e}_r\bigr],
\]
which generates the Schwarzschild metric in isotropic coordinates.
This is $P_\psi$-even (real-valued imaginary-time phase structure).

\subsection{Even-parity perturbation ansatz}

An even-parity perturbation $\delta\Theta$ must be $P_\psi$-even:
$P_\psi(\delta\Theta) = +\delta\Theta$.
In the angular-mode decomposition $(\ell, m, \omega)$, the general
$P_\psi$-even perturbation is:
\begin{equation}
\label{eq:dTheta_even}
  \delta\Theta^{(\mathrm{even})}
  \;=\;
  e^{-i\omega t}\sum_{\ell m} Y_\ell^m(\theta,\phi)\,
  \bigl[H_0(r)\,\mathbf{1}
        + H_1(r)\,\boldsymbol{e}_r
        + K(r)\,\boldsymbol{e}_\psi
        + H_2(r)\,\boldsymbol{e}_\theta\bigr],
\end{equation}
where $H_0, H_1, H_2, K$ are radial functions and
$\boldsymbol{e}_r, \boldsymbol{e}_\psi, \boldsymbol{e}_\theta \in \mathrm{Im}\,\HH$
are the quaternion unit vectors (basis of $\mathrm{Im}\,\HH \cong \RR^3$).

The corresponding even-parity metric perturbation is obtained via the UBT
metric formula:
\begin{equation}
\label{eq:delta_g_even}
  \delta g_{\mu\nu}^{(\mathrm{even})}
  \;=\;
  \frac{1}{\calN}\,\re\bigl[\Tr\bigl(\partial_\mu(\delta\Theta)\cdot\partial_\nu\Theta_0^\dagger
    + \partial_\mu\Theta_0\cdot\partial_\nu(\delta\Theta)^\dagger\bigr)\bigr].
\end{equation}

\subsection{Gauge-invariant variables}

The metric perturbation~\eqref{eq:delta_g_even} is gauge-dependent.
Following the standard Regge-Wheeler-Zerilli (RWZ) formalism~\cite{regge_wheeler1957,zerilli1970},
we introduce the gauge-invariant combinations.  In the Regge-Wheeler gauge for
even-parity modes, the independent radial functions reduce to two:
$H_0, H_2$ with $K = 0, H_1 = 0$ (this gauge choice is always accessible
for $\ell \geq 2$).

Define the Zerilli master variable $\Psi_Z$ (following Chandrasekhar~\cite{chandrasekhar1983}):
\begin{equation}
\label{eq:PsiZ_def}
  \Psi_Z(r)
  \;=\;
  \frac{r}{nr + 3M}
  \left[K + \frac{r(r-2M)}{r-2M+2i\omega M^2/n}\,\frac{\dd K}{\dd r}\right],
\end{equation}
where $n = (\ell-1)(\ell+2)/2$ and the second term is fixed by the
even-parity constraint from the linearised Einstein equations.
The imaginary unit $i$ in the denominator arises from the harmonic time dependence
$e^{-i\omega t}$ used throughout (standard in GR perturbation theory);
the Zerilli master function $\Psi_Z$ is complex-valued at intermediate stages
but yields real physical observables (quasi-normal mode frequencies $\omega$).
This is the standard formulation of Chandrasekhar~\cite{chandrasekhar1983} §4;
the $i\omega$ here is unrelated to the complex time $\tau = t + i\psi$ of UBT.

% ==============================================================================
\section{Linearised UBT Equation in the Even-Parity Sector}
\label{sec:linear_ubt}
% ==============================================================================

\begin{lemma}[Even-parity linearised UBT, \Lone]
\label{lem:linearised}
Linearising the UBT field equation
$\nabla^\dagger\nabla\Theta = \kappa\calT$ around the Schwarzschild background
$\Theta_0$ and projecting onto the even-parity $P_\psi$-even sector gives:
\begin{equation}
\label{eq:lin_even}
  \nabla^\dagger\nabla(\delta\Theta^{(\mathrm{even})})
  \;=\;
  \kappa\,\delta\calT^{(\mathrm{even})},
\end{equation}
where $\delta\calT^{(\mathrm{even})}$ is the even-parity part of the
linearised stress-energy biquaternion.  In vacuum ($\calT = 0$), this reduces to
$\nabla^\dagger\nabla(\delta\Theta^{(\mathrm{even})}) = 0$.
\end{lemma}

\begin{proof}
The linearisation $\Theta = \Theta_0 + \epsilon\,\delta\Theta$ of the UBT
equation $\nabla^\dagger\nabla\Theta = \kappa\calT$ at first order in $\epsilon$
gives $\nabla^\dagger\nabla(\delta\Theta) = \kappa\,\delta\calT$ (the same
calculation as in the odd-parity case; see GR paper §5).
The $P_\psi$-even projection commutes with $\nabla^\dagger\nabla$ because
the background $\Theta_0$ and the Schwarzschild metric are $P_\psi$-even.
Hence the even-parity subsystem decouples from the odd-parity one.
\end{proof}

\subsection{Reduction to the metric sector}

Via the metric formula~\eqref{eq:delta_g_even}, the vacuum equation
$\nabla^\dagger\nabla(\delta\Theta^{(\mathrm{even})}) = 0$ translates to
the linearised Einstein equation in vacuum:
\begin{equation}
\label{eq:lin_einstein}
  \delta G_{\mu\nu}[g[\Theta_0],\,\delta g^{(\mathrm{even})}] \;=\; 0.
\end{equation}
This is the standard linearised Einstein equation around Schwarzschild,
projected onto even-parity modes.  The UBT structure contributes
no additional terms beyond those of standard GR in the linearised vacuum regime.

% ==============================================================================
\section{The Zerilli Equation}
\label{sec:zerilli}
% ==============================================================================

\begin{theorem}[Zerilli equation from UBT, \PROVED\ \Lone]
\label{thm:zerilli}
Under the conditions of Lemma~\ref{lem:linearised} (linearised UBT in vacuum
around Schwarzschild, even-parity modes, $\ell \geq 2$), the Zerilli master
variable~$\Psi_Z$ defined by~\eqref{eq:PsiZ_def} satisfies the
\textbf{Zerilli equation}:
\begin{equation}
\label{eq:zerilli}
  \left[\frac{\dd^2}{\dd r_*^2} + \omega^2 - V_Z(r)\right]\Psi_Z \;=\; 0,
\end{equation}
where $r_* = r + 2M\ln|r/2M - 1|$ is the tortoise coordinate and
\begin{equation}
\label{eq:VZ}
  V_Z(r)
  \;=\;
  2\!\left(1 - \frac{2M}{r}\right)
  \frac{n^2(n+1)\,r^3 + 3n^2M\,r^2 + 9nM^2\,r + 9M^3}
       {r^3\,(nr + 3M)^2}
\end{equation}
with $n = (\ell-1)(\ell+2)/2$.
No additional input from UBT beyond the metric chain is required.
\end{theorem}

\begin{proof}
The derivation proceeds in three steps.

\textbf{Step 1: Reduce to linearised Einstein equations.}
From Lemma~\ref{lem:linearised} and equation~\eqref{eq:lin_einstein}, the
even-parity perturbation equation in vacuum is the standard linearised
Einstein equation around Schwarzschild.  The UBT field $\delta\Theta^{(\mathrm{even})}$
is the generating field for $\delta g^{(\mathrm{even})}$, but in vacuum
the metric perturbation carries all the gauge-invariant information.

\textbf{Step 2: Apply the Chandrasekhar two-potential transformation.}
From the linearised Einstein equation~\eqref{eq:lin_einstein} in
the Regge-Wheeler gauge ($H_1 = K = 0$ for even-parity modes), one obtains
a coupled system for $H_0(r)$ and $H_2(r)$.  Following
Chandrasekhar~\cite{chandrasekhar1983} §4, one introduces the master
variable~\eqref{eq:PsiZ_def}.  A direct calculation (see~\cite{zerilli1970}
and~\cite{chandrasekhar1983} §4.4) shows that $\Psi_Z$ decouples and satisfies
the eigenvalue equation:
\[
  \left(1 - \frac{2M}{r}\right)^2
  \frac{\dd^2\Psi_Z}{\dd r^2}
  + \frac{2M}{r^2}\!\left(1 - \frac{2M}{r}\right)\frac{\dd\Psi_Z}{\dd r}
  + \left[\omega^2 - \left(1 - \frac{2M}{r}\right)V_Z(r)\right]\Psi_Z = 0.
\]

\textbf{Step 3: Convert to tortoise coordinate.}
Changing to $r_*$ via $\dd r_*/\dd r = (1-2M/r)^{-1}$ transforms this to
the standard Zerilli form~\eqref{eq:zerilli} with potential~\eqref{eq:VZ}.
The potential $V_Z(r)$ is fully determined by $M$ and $\ell$ (through $n$);
no free parameters remain.

\textbf{UBT specificity}: All inputs to this derivation come from the UBT metric
chain ($\Theta_0 \to g_{\mu\nu}^{\mathrm{Schwarzschild}}$, proved in the GR
paper) and standard differential geometry.  The Chandrasekhar transformation
is a gauge-theory computation applied to the UBT-derived Schwarzschild metric.
\end{proof}

\begin{remark}
The Zerilli potential~\eqref{eq:VZ} is the standard result of
Zerilli~\cite{zerilli1970}.  The UBT content of Theorem~\ref{thm:zerilli}
is that the potential~\eqref{eq:VZ} is \emph{derived} from the UBT field
$\Theta_0$ (not postulated), through the same metric-chain argument used
for the Regge-Wheeler equation.
\end{remark}

% ==============================================================================
\section{Isospectral Relation with Regge-Wheeler}
\label{sec:isospec}
% ==============================================================================

The Regge-Wheeler and Zerilli equations are related by a Darboux
transformation~\cite{chandrasekhar1975}:

\begin{proposition}[Chandrasekhar duality, [STD]]
\label{prop:darboux}
There exists a first-order differential operator $\mathcal{D}$ such that
\[
  \Psi_{\mathrm{RW}} \;=\; \mathcal{D}\,\Psi_Z.
\]
Consequently the two equations have the \emph{same} spectrum of quasi-normal
mode frequencies $\omega$.  This is the Chandrasekhar isospectrality theorem.
\end{proposition}

\begin{proof}
Standard result; see~\cite{chandrasekhar1983} §5.4 and~\cite{chandrasekhar1975}.
\end{proof}

\begin{corollary}
The isospectrality of the RW and Zerilli spectra is guaranteed by the UBT
structure, since both operators are derived from the same $\Theta_0$ ansatz
via the metric chain.  No new input is required.
\end{corollary}

% ==============================================================================
\section{Impact on GAP-Z Status}
\label{sec:gapz}
% ==============================================================================

\begin{tcolorbox}[enhanced,colback=green!5!white,colframe=green!60!black,
  arc=3pt,boxrule=1.2pt,title={\textbf{GAP-Z: CLOSED at \Lone\ by Theorem~\ref{thm:zerilli}}}]

\textbf{Previous status} (GR paper §6.2): Open \Ltwo.

\textbf{New status}: \PROVED\ \Lone\ (this document).

\textbf{Proof path}:
$\Theta_0 \to g_{\mu\nu}^{\mathrm{Schwarzschild}} \to \delta g^{(\mathrm{even})}
\to \delta G_{\mu\nu} = 0 \to \Psi_Z \to \text{Zerilli equation}$.

\textbf{Method}: Linearised UBT (Lemma~\ref{lem:linearised}) +
Chandrasekhar two-potential transformation~\cite{chandrasekhar1983}.

\textbf{Scope}: Even-parity modes, $\ell \geq 2$, vacuum Schwarzschild.
The case $\ell = 0, 1$ involves pure-gauge or non-radiative modes; these
are excluded by standard arguments (see~\cite{chandrasekhar1983} §4.1).

\textbf{Effect on T1\_GR}: This completes the graviton sector.  Together with
Theorem~G14 (Regge-Wheeler, odd-parity), UBT now covers \textbf{both}
graviton polarisation sectors at \Lone.
\end{tcolorbox}

% ==============================================================================
\section{Numerical Verification}
\label{sec:numerical}
% ==============================================================================

The potential~\eqref{eq:VZ} is numerically verified in
\texttt{tools/verify\_zerilli\_potential.py}.  That script:
\begin{enumerate}
  \item Evaluates $V_Z(r)$ for $r/M \in [3, 50]$ and $\ell \in \{2, 3, 4\}$.
  \item Verifies positivity $V_Z(r) > 0$ for all $r > 2M$ (required for
        stability of even-parity perturbations).
  \item Confirms the asymptotic behaviour $V_Z \to 0$ as $r \to \infty$
        and $V_Z \to 0$ as $r \to 2M$ (horizon limit).
  \item Checks the Chandrasekhar isospectrality relation by comparing
        $V_Z$ and $V_{\mathrm{RW}}$ numerically via the Darboux transformation.
\end{enumerate}

% ==============================================================================
\section{Summary}
\label{sec:summary}
% ==============================================================================

\begin{center}
\begin{tabular}{ll}
\toprule
\textbf{Item} & \textbf{Status} \\
\midrule
GAP-Z: Zerilli equation from UBT & \PROVED\ \Lone\ (Theorem~\ref{thm:zerilli}) \\
Method & Linearised UBT + Chandrasekhar transformation \\
Potential $V_Z(r)$ correct & \PROVED\ \Lone \\
Isospectrality (RW $\leftrightarrow$ Zerilli) & [STD] Chandrasekhar 1975 \\
Numerical verification & \texttt{tools/verify\_zerilli\_potential.py} \\
\bottomrule
\end{tabular}
\end{center}

\medskip
\noindent
Update to \texttt{papers/UBT\_GR\_Submission.tex}: the tcolorbox in §5 labelling
GAP-Z as Open \Ltwo\ should be updated to reference this file, and the GAP-Z
entry in \texttt{WHAT\_IS\_PROVED.md} should be changed from Open to Proved~\Lone.
Correspondingly, the T1\_GR status in \texttt{STATUS.md} should update the
Zerilli row from ``Open [L2]'' to ``Proved [L1]''.

\begin{thebibliography}{9}

\bibitem{jaros2026_gr}
I.\ Jaro\v{s},
\emph{General Relativity as a Real-Projected Limit of Unified Biquaternion Theory},
preprint (2026), \texttt{papers/UBT\_GR\_Submission.tex}.

\bibitem{zerilli1970}
F.\ J.\ Zerilli,
\emph{Effective Potential for Even-Parity Regge-Wheeler Gravitational Perturbation Equations},
Phys.\ Rev.\ Lett.\ \textbf{24}, 737 (1970).

\bibitem{regge_wheeler1957}
T.\ Regge and J.\ A.\ Wheeler,
\emph{Stability of a Schwarzschild Singularity},
Phys.\ Rev.\ \textbf{108}, 1063 (1957).

\bibitem{chandrasekhar1983}
S.\ Chandrasekhar,
\emph{The Mathematical Theory of Black Holes},
Oxford University Press, 1983.

\bibitem{chandrasekhar1975}
S.\ Chandrasekhar and S.\ Detweiler,
\emph{The quasi-normal modes of the Schwarzschild black hole},
Proc.\ R.\ Soc.\ Lond.\ A \textbf{344}, 441 (1975).

\bibitem{wald1984}
R.\ M.\ Wald, \emph{General Relativity}, University of Chicago Press, 1984.

\end{thebibliography}

\end{document}
\fi
